\section{问题引入：为什么“接上知识库”仍然不够}
\label{sec:intro}

大语言模型将大量知识存储在参数中，训练完成后再更新这些知识并不容易。模型对热门实体往往回答得很好，遇到长尾事实、近期事件或专业领域细节时，却可能给出流畅而缺乏依据的内容。检索增强生成把外部文档作为推理时可访问的非参数记忆，使知识库更新和答案溯源更加灵活。REALM 在预训练阶段引入检索，Lewis 等提出通用的 RAG 序列生成框架，稠密段落检索则为开放域问答提供了可学习的文档访问方式~\citep{guu2020realm,lewis2020rag,karpukhin2020dense}。

常见的 RAG 实现先对用户问题编码，取回固定数量的文档，再把文档与问题一起交给生成器。这一设计简单有效，却也带来了固定的行为模式。以“水的化学式是什么”为例，模型已有知识通常足够，额外检索主要增加延迟；询问“某部获奖电影导演的生日”时，系统需要先确定导演，才能形成下一步查询，一次检索很难覆盖全部证据。Mallen 等的实验也呈现出类似差异：不借助检索的模型在高流行度实体问题上仍有竞争力，检索则在长尾知识上更有价值~\citep{mallen2023trust}。因此，检索是否值得执行，需要结合问题和模型判断。

文档进入上下文以后，仍有一段距离才能转化为可靠答案。长上下文研究发现，关键证据位于开头或结尾时通常更容易被利用，放在中部则可能明显降低表现~\citep{liu2024lost}。Cuconasu 等考察了不同噪声文档的作用，显示文档数量、相关程度和混合方式都会改变生成结果~\citep{cuconasu2024noise}。检索器找到了相关段落，生成器仍可能忽略其中的限制条件，或者被同名实体和过期信息带偏。检索、上下文利用和答案忠实性因而需要分别讨论。

已有综述通常以朴素 RAG、高级 RAG 和模块化 RAG 描述技术演进，并围绕检索、生成和增强组件组织内容~\citep{gao2024survey}。模块化 RAG 进一步引入路由、循环和融合等结构，使复杂系统可以拆成可组合的算子~\citep{gao2024modular}。在这些分类基础上，本文把视线放在系统中的决策行为：什么信号促使模型检索，查询如何随中间结果变化，证据不足时采取什么行动，以及不断增加的检索在什么条件下失去价值。

文章先建立统一的状态与行动表示，再按检索触发、查询构建、证据评价和停止控制梳理代表方法，随后讨论迭代推理与评价方案。最后的个人分析围绕两个具体问题展开：怎样记录当前还缺哪些证据，以及怎样判断下一次检索是否仍值得执行。

